\scanpage{118}
% archive foliation: 117
III) a) Si $s,t\in\Gamma(U,\mathcal P_X)$, $U$ connexe, \deleted{\unsure{donc}} et $G(t,s)=0$, alors $\exists\deleted{\unsure{unique}}\ g\in H$ tel que $t=g\cdot s$.

b) Pour $t$ fixé, le morphisme
\[
s\longmapsto G(t,s)\qquad \mathcal P_X\longrightarrow\underline\omega_X^{\otimes2}
\]
(qui induit un hom. injectif $\deleted{\unsure{\mathcal P_X/H}}\ H\backslash\mathcal P_X\lhook\joinrel\longrightarrow\underline\omega_X^{\otimes2}$, d'après a)) est \uline{surjectif} (i.e.
\[
H\backslash\mathcal P_X\xrightarrow{\sim}\underline\omega_X^{\otimes2}\qquad s\longmapsto G(t,s)=\left(-2\frac{s'''}{s'}+3\left(\frac{s''}{s'}\right)^2\right)dt^2
\]

D'autre part, grâce à II, on peut donner une forme plus simple et intrinsèque à l'expression $G(t,s)$.

\uline{Proposition} (qui est formelle à partir de (II)) On peut, \inserted{à isomorphisme unique près,} trouver un torseur $\operatorname{Conn}_X$ sous $\underline\omega_X^{\otimes2}$, et un morphisme de faisceaux
\[
t\longmapsto c_t:\quad\mathcal P_X\xrightarrow{c^X}\operatorname{Conn}_X
\]
tels que l'on ait, pour $s,t\in\Gamma(U,\mathcal P_X)$
\[
\boxed{G(t,s)=c_s-c_t}
\]

\margin{!}
Le faisceau $\operatorname{Conn}_X$ est appelé le \uline{faisceau des connexions projectives} (analytiques complexes) sur $X$. Une section dudit est appelé \uline{connexion projective} (analytique) de $X$. Le morphisme $c^X$ est inv. par $H$ et induit par passage au quotient un iso. can.
\[
c^X:H\backslash\mathcal P_X\xrightarrow{\sim}\operatorname{Conn}_X.
\]
Un morphisme étale
\[
f:X'\longrightarrow X
\]
définit des iso.
\[
\mathcal P_{X'}\simeq f^{-1}(\mathcal P_X)\qquad \operatorname{Conn}_{X'}\simeq f^{-1}(\operatorname{Conn}_X)
\]
et on trouve un \deleted{\unsure{même}} diag. comm.
\diagram{\begin{tikzcd}
\mathcal P_{X'}\arrow[r,"c^{X'}"]\arrow[d,"\sim"']&\deleted{\unsure{\mathcal C_{X'}}}\operatorname{Conn}_{X'}\arrow[d,"\sim"]\\
f^{-1}(\mathcal P_{\unsure{X'|X}})\arrow[r,"c^X"']&f^{-1}(\operatorname{Conn}_X)
\end{tikzcd}}
